% from ega1-7-fr.tex
\oldpage[I]{157}
tout ouvert partout dense $U$ de $X$, $U_i=U\cap X'_i$ est un ouvert non vide
dense dans $X'_i$, les $U_i$ étant deux à deux sans point commun, donc
$U'=\bigcup_i U_i$ est dense dans $X$. Se donner un morphisme de $U'$ dans $Y$
revient à se donner (arbitrairement) un morphisme de chacun des $U_i$ dans $Y$.
Donc~:
\end{env}

\begin{proposition}[7.1.7]
\label{I.7.1.7-fr}
Soient $X$, $Y$ deux préschémas (resp. $S$-préschémas) tels que $X$ ait un
nombre fini de composantes irréductibles $X_i$, de points génériques $x_i$
($1\leq i\leq n$). Si $R_i$ est l'ensemble des germes de morphismes (resp.
$S$-morphismes) de parties ouvertes de $X$ dans $Y$ au point $x_i$, l'ensemble
des applications rationnelles (resp. $S$-applications rationnelles) de $X$
dans $Y$ s'identifie au produit des $R_i$ ($1\leq i\leq n$).
\end{proposition}

\begin{corollary}[7.1.8]
\label{I.7.1.8-fr}
Soit $X$ un préschéma noethérien. L'anneau des fonctions rationnelles sur $X$
est un anneau artinien, dont les composants locaux sont les anneaux
$\mathscr{O}_{x_i}$ des points génériques $x_i$ des composantes irréductibles
de $X$.
\end{corollary}

\begin{corollary}[7.1.9]
\label{I.7.1.9-fr}
Soient $A$ un anneau noethérien, et soit $X=\operatorname{Spec}(A)$. Si $Q$ est
le complémentaire de la réunion des idéaux premiers minimaux de $A$, l'anneau
des fonctions rationnelles sur $X$ s'identifie canoniquement à l'anneau de
fractions $Q^{-1}A$.
\end{corollary}

Cela résultera du lemme suivant~:

\begin{lemma}[7.1.9.1]
\label{I.7.1.9.1-fr}
Pour qu'un élément $f\in A$ soit tel que $D(f)$ soit partout dense dans $X$, il
faut et il suffit que $f\in Q$~; tout ouvert dense dans $X$ contient un ouvert
de la forme $D(f)$, où $f\in Q$.
\end{lemma}

En effet, supposons ce lemme démontré~; l'anneau des sections
$\Gamma(D(f),\mathscr{O}_X)$ s'identifiant à $A_f$
(\hyperref[I.1.3.6-fr]{1.3.6} et \hyperref[I.1.3.7-fr]{1.3.7}), il résulte du
fait que les $D(f)$ avec $f\in Q$ forment un ensemble cofinal dans l'ensemble
ordonné (pour $\supset$) des ouverts denses dans $X$, et de la déf.
(\hyperref[I.7.1.1-fr]{7.1.1}) que l'anneau des fonctions rationnelles sur $X$
s'identifie à la limite inductive des $A_f$ pour $f\in Q$ (pour la relation de
préordre «~$g$ est multiple de $f$~»), c'est-à-dire à $Q^{-1}A$
(\hyperref[0.1.4.5-fr]{0, 1.4.5}).

Pour démontrer (\hyperref[I.7.1.9.1-fr]{7.1.9.1}), désignons encore par $X_i$
les composantes irréductibles de $X$ ($1\leq i\leq n$)~; si $D(f)$ est dense
dans $X$, $D(f)\cap X_i\neq\varnothing$ pour $1\leq i\leq n$ et
réciproquement~; mais cela signifie que $f\notin\mathfrak{p}_i$ pour
$1\leq i\leq n$, en posant $\mathfrak{p}_i=\mathfrak{j}(X_i)$, et comme les
$\mathfrak{p}_i$ sont les idéaux premiers minimaux de $A$
(\hyperref[I.1.1.14-fr]{1.1.14}), les relations
$f\notin\mathfrak{p}_i$ ($1\leq i\leq n$) équivalent à $f\in Q$, d'où la
première assertion du lemme. D'autre part, si $U$ est un ouvert dense dans $X$,
le complémentaire de $U$ est un ensemble de la forme $V(\mathfrak{a})$, où
$\mathfrak{a}$ est un idéal qui n'est contenu dans aucun des $\mathfrak{p}_i$~;
il n'est donc pas contenu dans leur réunion (\cite{I-10}, p.~13), et il existe
donc un $f\in\mathfrak{a}$ appartenant à $Q$~; d'où $D(f)\subset U$, ce qui
achève la démonstration.

\begin{env}[7.1.10]
\label{I.7.1.10-fr}
Supposons de nouveau $X$ irréductible, de point générique $x$. Comme tout
ouvert non vide $U$ de $X$ contient $x$, et par suite contient aussi tout
$z\in X$ tel que $x\in\overline{\{z\}}$, tout morphisme $U\to Y$ peut se
composer avec le morphisme canonique $\operatorname{Spec}(\mathscr{O}_x)\to X$
(\hyperref[I.2.4.1-fr]{2.4.1})~; et deux morphismes dans $Y$ de deux parties
ouvertes non vides de $X$, qui coïncident dans une partie ouverte non vide de
$X$ donnent par composition le même morphisme
$\operatorname{Spec}(\mathscr{O}_x)\to Y$. Autrement dit, à toute application
rationnelle de $X$ dans $Y$ correspond ainsi un morphisme
$\operatorname{Spec}(\mathscr{O}_x)\to Y$ bien déterminé.
\end{env}

